% !TEX root = ../main.tex
\subsection{问题三：固定尺寸下的全域干燥终点}
\subsubsection{后期环境延长与终点数学定义}
问题三沿用式\eqref{eq:q2pde}，半径仍为$R_0=0.02$ m。附件1仅覆盖4 h，但模型预期运行两三天，因此必须对数据之外的边界输入进行闭合。观察后期输入已接近稳定，取3--4 h共61个观测点的算术均值：
\begin{equation}
T_a(t)=49.9989344262\ ^\circ\mathrm C,\qquad
C_e(t)=0.0499875410,\qquad t>14400\ \mathrm s .
\label{eq:plateau}
\end{equation}
该方法不把最后单个测量值的波动永久延伸到全部后期。4 h前保持线性插值；4 h处分段结束积分，以当时内部状态作为后一段初值，避免跨过输入规则的微小跳变。内部温湿状态连续，外部输入可以在分段点按指定规则切换。

对于“药材各处的水分浓度均低于0.15 kg/kg”，定义
\begin{equation}
M(t)=\max_{0\le r\le R_0}C(r,t),\qquad
g(t)=M(t)-0.15 .
\end{equation}
干燥要求等价于$g(t)<0$。对于本次连续下降并穿越阈值的轨迹，定义临界时间
\begin{equation}
t_*=\inf\{t\ge0:M(t)<0.15\},\qquad M(t_*)=0.15 .
\label{eq:event}
\end{equation}
严格小于通常不存在包含边界的“最早等号时刻”，因此本文报告的时长是达标集合的下确界，即穿越临界。实际要求严格小于时，应略晚于该时刻停止。

采用最大值而不采用平均值，有明确的空间含义。平均含水率
\begin{equation}
\overline C(t)=
\frac{\int_V\rho_d C\dd V}{\int_V\rho_d\dd V}
=\frac{2}{R_0^2}\int_0^{R_0}C(r,t)r\dd r
\end{equation}
只描述总体失水程度。即使外层已经很干，只要中心仍超过阈值，就不满足题目“各处”的要求。相反，仅以表面为准容易过早停止，因为表面直接与环境交换，通常最早接近低含水状态。

\note{后期环境延长是与数据范围有关的物理假设，不能靠细化空间网格消除。第六部分利用既有末值保持结果报告该假设的敏感性。}

\nextpage
\subsubsection{连续事件定位与长时间过程求解}
计算中在每次事件函数评估时检查全部径向节点，使用
\begin{equation}
g_N(t)=\max_{0\le i\le N}C_i(t)-0.15 .
\end{equation}
程序设置事件下降穿越方向，并在变号区间内使用求解器的连续数值表示定位零点。每60 s导出一次只是结果采样规则，不是终点定位精度。若直接等到下一条60 s采样记录才宣布完成，便会把输出间隔误当成计算误差。

在本次结果中，全域最大值始终由中心节点控制，最终最大值位置为$x=0$。这与表面先失水的径向分布一致，但程序没有把“中心必然最后干”写成先验简化，而是保留了对全体节点的检查。空间网格与保形插值使输出曲线不会凭空制造明显的内部超调，连续场近似误差再由加密比较衡量。

\fig{fig8}{0.98}{已计算的最大含水率演化与半径变化；两种方案物性也不同，不能只归因为收缩}{fig:drying}

图\ref{fig:drying}左侧显示，初期最大含水率下降相对较慢，随后进入明显失水阶段，接近阈值时又出现长尾。早期中心尚未充分响应外层交换；中期温度提高促进扩散；后期低含水率使扩散系数降低。三者共同造成干燥曲线并非固定斜率，也难以用单一指数在全时段精确替代。

主模型的临界时刻为
\begin{equation}
t_*^{(3)}=206906.3987685\ \mathrm s
=57.4739996579\ \mathrm h .
\end{equation}
按题目四位小数报告为\textbf{57.4740 h}。若按整秒给出严格低于的停止时刻，可取206907 s；该区别来自严格不等式与数值报告规则，不是新的物理干燥机制。

\nextpage
\subsubsection{全程水分分布与终点结果}
\begin{table}[H]\centering
\caption{问题3：固定半径下全程干基含水率（kg/kg）}\label{tab:q3}
\begin{tabular}{cccccc}\toprule
时间/h & 0 cm & 0.5 cm & 1 cm & 1.5 cm & 2 cm\\ \midrule
6.0000 & 1.0170 & 0.9881 & 0.9015 & 0.7550 & 0.5328\\
12.0000 & 0.4561 & 0.4432 & 0.4031 & 0.3297 & 0.1642\\
18.0000 & 0.2988 & 0.2916 & 0.2687 & 0.2253 & 0.0845\\
24.0000 & 0.2378 & 0.2327 & 0.2167 & 0.1855 & 0.0665\\
30.0000 & 0.2058 & 0.2018 & 0.1892 & 0.1641 & 0.0598\\
36.0000 & 0.1859 & 0.1825 & 0.1718 & 0.1504 & 0.0566\\
42.0000 & 0.1721 & 0.1691 & 0.1597 & 0.1407 & 0.0548\\
48.0000 & 0.1618 & 0.1592 & 0.1507 & 0.1334 & 0.0537\\
54.0000 & 0.1539 & 0.1514 & 0.1436 & 0.1277 & 0.0530\\
57.4740 & 0.1500 & 0.1477 & 0.1402 & 0.1249 & 0.0526\\
\bottomrule\end{tabular}\end{table}


表\ref{tab:q3}给出每6 h的五个径向位置及临界终点。12 h时中心含水率为0.4561 kg/kg，表面为0.1642 kg/kg；18 h时表面已低至0.0845 kg/kg，而中心仍为0.2988 kg/kg。由此可直接看出，以表面含水率决定停止会遗漏仍偏湿的内部区域。

48 h时，距中心1 cm处含水率为0.1507 kg/kg，已接近阈值，但中心为0.1618 kg/kg。54 h时中心仍为0.1539 kg/kg，因此干燥还需持续约3.4740 h。临界终点处，中心为0.1500，其他四个展示位置均低于0.15，表面已降至0.0526 kg/kg。

临界表值保留四位小数，不能用显示为0.1500判断原始值严格大于、等于或小于阈值。连续事件定位使用未舍入的状态；结果表在计算完成后才舍入。该顺序保证“判断精度”不被“显示精度”替代。

既有\texttt{result3.xlsx}按60 s间隔保存各位置含水率，并增加临界时刻末行。每个时间点对应0--2 cm的21个位置，满足正文表格之外的空间明细需求。末行不强行对齐60 s整数倍，因为其任务是给出连续模型定位的终点状态。

从工艺含义看，后期外层接近环境平衡而中心缓慢下降，继续提高外部传质能力未必能同比缩短全部干燥时间；内部扩散逐渐成为重要限制。本文没有实施额外工艺参数优化，因此此处仅依据现有分布解释潜在限制，不报告未经计算的节能比例或最优控制策略。
